\documentclass[11pt, a4paper]{article}

% ── Packages ──────────────────────────────────────────────────────────────────
\usepackage{fontspec}
\usepackage{amsmath, amssymb}
\usepackage{graphicx}
\usepackage{booktabs}
\usepackage{hyperref}
\usepackage[margin=1in]{geometry}
\usepackage{natbib}
\usepackage{xcolor}
\usepackage{enumitem}
\usepackage{float}

\hypersetup{
    colorlinks=true,
    linkcolor=blue!60!black,
    citecolor=blue!60!black,
    urlcolor=blue!60!black,
}

% ── Hieroglyphic font ─────────────────────────────────────────────────────────
\newfontfamily\hierofont{EgyptianHiero}[
    Path=../../final_output/,
    Extension=.ttf,
]
\newcommand{\hiero}[1]{{\normalfont\hierofont #1}}

% ── Title ─────────────────────────────────────────────────────────────────────
\title{
    \textbf{Heiroglyphy: Recovering the Conceptual Geometry\\of Ancient Egyptian Through Vector Space Alignment}
}

\author{
    Erik Brinsmead \\
    Independent Researcher \\[6pt]
    Claude \\
    Anthropic \\[12pt]
    \texttt{github.com/ebrinz/heiroglyphy}
}

\date{March 2026}

\begin{document}
\maketitle

% ══════════════════════════════════════════════════════════════════════════════
\begin{abstract}
Ancient Egyptian has been translated by scholars for centuries, but translation is inherently lossy. When \textit{n\d{t}r} becomes ``god,'' the web of relationships that defined divinity for an Egyptian speaker---its proximity to kingship, offering, truth, and cosmic order---is compressed into a single modern word. We present a method for recovering this lost semantic structure using vector space alignment. By training word embeddings on 100,729 Ancient Egyptian sentences and aligning the resulting geometric space to English, we construct a bridge that preserves not just word meanings but the \textit{relationships between meanings}. Across 15 experimental iterations, we achieve 32.35\% Top-1 unsupervised cross-lingual retrieval accuracy---without a bilingual dictionary---and demonstrate that the aligned space reveals conceptual structures invisible to literal translation: that gold and divinity occupy the same semantic region, that silence is the defining condition of death, that the Egyptian serpent is divine rather than wise, and that the ``earth mother'' archetype is absent from the Egyptian conceptual landscape. The methodology is domain-agnostic: anywhere two signal spaces share latent structure, this approach can recover it.
\end{abstract}

% ══════════════════════════════════════════════════════════════════════════════
\section{Introduction}
\label{sec:intro}

Imagine two cities, each with their own map. The streets are labeled in different languages, the neighborhoods have different names, and the scale is different. But the \textit{shape} of the city---the way neighborhoods cluster together, the distance between the market and the harbor, the fact that the temple sits between the palace and the cemetery---is the same. If you could find the right rotation and scaling, you could overlay one map onto the other, and every landmark would align.

This is the core hypothesis of this work: that languages, when represented as geometric objects, share a common shape. Words that mean similar things cluster together in every language. ``Water'' sits near ``river'' and ``flood'' whether you write it in English or in Middle Egyptian. The challenge is to find the transformation that aligns one language's shape onto another's.

What makes this interesting for Ancient Egyptian is not the translation itself---Egyptologists have been translating hieroglyphic texts for two centuries. It is what the \textit{geometry} reveals that translation cannot. A dictionary tells you that \textit{n\d{t}r} means ``god.'' But the embedding space tells you that \textit{n\d{t}r} sits near \textit{nbw} (gold), near \textit{nswt} (king), near \textit{m\=a\c{a}.t} (truth/cosmic order)---and \textit{far} from where an English speaker would place ``god'' relative to ``gold.'' The distance between ``god'' and ``gold'' in Egyptian is nearly zero. In English, it is vast. This difference is invisible to translation. It is visible to geometry.

The purpose of this project is to build that geometric lens---to align the embedding space of Ancient Egyptian to English with enough fidelity that the \textit{relational structure} of the ancient language becomes readable. The accuracy metric (32.35\% Top-1 retrieval) is a means to an end. The end is the discoveries: the conceptual architecture of a civilization that thought differently from us, preserved in the statistical structure of its texts.

% ══════════════════════════════════════════════════════════════════════════════
\section{Background: Words as Geometry}
\label{sec:background}

\subsection{The Distributional Hypothesis}

The foundation of this work is a simple observation from linguistics: words that appear in similar contexts tend to have similar meanings \citep{harris1954distributional, firth1957synopsis}. ``Doctor'' and ``physician'' appear near the same words (``patient,'' ``hospital,'' ``treat''), so they must mean similar things. This is the \textit{distributional hypothesis}.

In the 2010s, researchers discovered that this hypothesis could be operationalized at scale \citep{mikolov2013efficient}. By training a neural network to predict a word from its context (or vice versa), each word in a language can be assigned a vector---a point in high-dimensional space---such that words with similar meanings are close together. The resulting space has remarkable properties: vector arithmetic encodes semantic relationships. The vector for ``king'' minus ``man'' plus ``woman'' yields a vector near ``queen'' \citep{mikolov2013linguistic}.

\subsection{Cross-Lingual Alignment}

A natural question follows: if every language forms a geometric shape, do different languages form the \textit{same} shape? Research on cross-lingual word embeddings suggests they do \citep{conneau2018word, lample2018word}. The word for ``water'' clusters near the word for ``river'' in every language. The overall topology is approximately isomorphic.

This means that if we can find a transformation---a rotation, scaling, and translation---that maps one language's embedding space onto another's, we can translate between them without a bilingual dictionary. The transformation can be learned from a small set of ``anchor'' pairs (words whose translation is known) and then applied to the entire vocabulary.

The most common approach is \textit{Procrustes alignment} \citep{schonemann1966generalized}: given anchor pairs $(x_i, y_i)$ where $x_i$ is a source-language vector and $y_i$ is a target-language vector, find the linear transformation $W$ that minimizes:

\begin{equation}
    W^* = \arg\min_W \sum_{i} \| Wx_i - y_i \|^2
\end{equation}

This has a closed-form solution via singular value decomposition (SVD) when $W$ is constrained to be orthogonal, or via Ridge regression when it is not.

\subsection{The Ancient Egyptian Challenge}

Applying cross-lingual alignment to Ancient Egyptian is harder than aligning, say, French and English:

\begin{itemize}[nosep]
    \item \textbf{No living speakers.} There is no bilingual community to validate results.
    \item \textbf{Sparse data.} The entire digitized corpus is $\sim$100,000 sentences, with a median length of 6 words. Modern NLP models train on billions of sentences.
    \item \textbf{Domain bias.} The surviving texts are overwhelmingly funerary and religious. Everyday speech is almost entirely lost.
    \item \textbf{Morphological complexity.} Egyptian words are agglutinative, with suffix pronouns, plural markers, and editorial annotations embedded in the transliteration.
    \item \textbf{Semantic distance.} Ancient Egyptian is 4,000 years removed from Modern English. The conceptual categories are different. There is no guarantee that the embedding spaces are isomorphic.
\end{itemize}

Despite these challenges, the distributional hypothesis does not require a living language. It requires only a corpus large enough to estimate co-occurrence statistics. The BBAW hieroglyphic text database \citep{bbaw_tla} provides exactly that.

% ══════════════════════════════════════════════════════════════════════════════
\section{Methodology}
\label{sec:methodology}

\subsection{Overview}

The final pipeline, arrived at through 15 experimental iterations, consists of four stages:

\begin{enumerate}[nosep]
    \item Train 768-dimensional FastText embeddings on the Egyptian transliteration corpus.
    \item Fuse text embeddings with 768-dimensional visual features (zero-padded) to create 1536-dimensional representations.
    \item Learn a Ridge regression mapping from the fused Egyptian space to the 300-dimensional English GloVe space, using 5,360 anchor pairs for training.
    \item Evaluate via nearest-neighbor retrieval on 1,340 held-out anchor pairs.
\end{enumerate}

\subsection{Data}

\textbf{Egyptian corpus.} 100,729 sentences of transliterated hieroglyphic text from the Berlin-Brandenburg Academy's Thesaurus Linguae Aegyptiae (TLA) \citep{bbaw_tla}, accessed via the BBAW HuggingFace dataset. Total tokens: 789,159. Unique vocabulary: 80,662 (before filtering).

\textbf{Anchor pairs.} 8,541 Egyptian--English word pairs extracted via co-occurrence analysis from parallel German translations in the TLA corpus. Each pair has a confidence score (mean: 0.747, range: 0.30--1.00). After filtering for vocabulary coverage, 6,700 pairs are used (80/20 train/test split, \texttt{random\_state=42}).

\textbf{English embeddings.} GloVe 6B 300d \citep{pennington2014glove}: 400,000 English words trained on 6 billion tokens of Wikipedia and Gigaword.

\textbf{Visual features.} 768-dimensional ResNet-50 features extracted from the HamdiJr/Egyptian\_hieroglyphs dataset \citep{hamdijr_hieroglyphs}: 4,210 labeled hieroglyph images spanning 115 Gardiner sign classes. Visual match rate in the final model: 0.08\%.

\subsection{Egyptian Embeddings}

We train FastText skip-gram embeddings \citep{bojanowski2017enriching} on the Egyptian transliteration corpus. FastText is preferred over Word2Vec because its subword (character n-gram) architecture captures morphological structure---critical for Egyptian's agglutinative morphology, where variants like \textit{n\d{t}r}, \textit{n\d{t}rw}, and \textit{n\d{t}r.j} share a common root.

The final training configuration, selected through a 9-configuration parameter sweep (Section~\ref{sec:journey}):

\begin{table}[H]
\centering
\begin{tabular}{ll}
\toprule
Parameter & Value \\
\midrule
\texttt{vector\_size} & 768 \\
\texttt{window} & 10 \\
\texttt{min\_count} & 5 \\
\texttt{sg} & 1 (skip-gram) \\
\texttt{epochs} & 10 \\
\bottomrule
\end{tabular}
\caption{FastText training parameters (V15 configuration).}
\label{tab:fasttext_params}
\end{table}

Two parameters were critical discoveries:

\textbf{\texttt{min\_count=5}} filters words appearing fewer than 5 times. The original model (\texttt{min\_count=1}) retained 80,662 words, of which 52,813 (65.5\%) were hapax legomena---words appearing exactly once. These single-occurrence words have essentially random vectors. Filtering to \texttt{min\_count=5} reduces vocabulary to 10,833 words (an 87\% reduction) while \textit{improving} alignment accuracy. Less noise, better geometry.

\textbf{\texttt{window=10}} doubles the context window from the original 5. The Egyptian corpus has a median sentence length of 6 words. With \texttt{window=5}, most words see only 2--3 context neighbors. Widening to 10 allows FastText to capture relationships across the full span of typical sentences.

\subsection{Fusion and Alignment}

Each Egyptian word receives a 1536-dimensional fused representation:

\begin{equation}
    \mathbf{v}_{\text{fused}} = [\mathbf{v}_{\text{text}} \| \mathbf{v}_{\text{visual}}]
\end{equation}

where $\mathbf{v}_{\text{text}} \in \mathbb{R}^{768}$ is the FastText vector and $\mathbf{v}_{\text{visual}} \in \mathbb{R}^{768}$ is the ResNet-50 visual feature vector (or $\mathbf{0}$ if no visual mapping exists). In practice, 99.92\% of words have zero visual features. The 1536-dimensional space nonetheless outperforms the 768-dimensional text-only space, likely due to implicit regularization from the zero-padding (Section~\ref{sec:journey}).

The alignment is learned via Ridge regression:

\begin{equation}
    W^* = \arg\min_W \| XW - Y \|_F^2 + \alpha \| W \|_F^2
\end{equation}

where $X \in \mathbb{R}^{n \times 1536}$ contains fused Egyptian vectors, $Y \in \mathbb{R}^{n \times 300}$ contains target GloVe vectors, and $\alpha$ is the regularization strength. The optimal $\alpha = 0.001$ was found through grid search evaluated by retrieval accuracy (not MSE---a critical distinction discussed in Section~\ref{sec:journey}).

All projected Egyptian vectors are L2-normalized before retrieval.

\subsection{Evaluation}

For each test pair $(e_i, g_i)$, we project the Egyptian vector $\mathbf{v}_{e_i}$ to English space and retrieve its $k$ nearest neighbors by cosine similarity in the full 400,000-word GloVe vocabulary. The pair is ``correct'' if $g_i$ appears in the top-$k$ results.

\begin{table}[H]
\centering
\begin{tabular}{lr}
\toprule
Metric & Score \\
\midrule
Top-1 Accuracy & 32.35\% \\
Top-5 Accuracy & 41.47\% \\
Top-10 Accuracy & 45.13\% \\
\midrule
Test pairs & 1,340 \\
Anchor coverage & 78.4\% \\
Visual match rate & 0.08\% \\
\bottomrule
\end{tabular}
\caption{Final evaluation metrics (V15).}
\label{tab:final_results}
\end{table}

% ══════════════════════════════════════════════════════════════════════════════
\section{The Iterative Journey}
\label{sec:journey}

The final methodology emerged through 15 experimental iterations. Each failure was as informative as each success.

\begin{table}[H]
\centering
\small
\begin{tabular}{clrc}
\toprule
V & Technique & Top-1 & Status \\
\midrule
1 & Neural Vec2Vec (Multi-Space) & --- & Failed \\
2 & Unsupervised Adversarial & --- & Failed \\
3 & Orthogonal Procrustes & 22.0\% & Baseline \\
4 & Procrustes + CSLS & 15.0\% & Regression \\
5 & Procrustes + 10$\times$ Data & 24.5\% & Improved \\
6 & BERT Contextual & 0.5\% & Failed \\
7 & FastText 768d & 29.1\% & Breakthrough \\
8 & Coptic Bridge & 28.2\% & Regression \\
9 & Text + Visual Fusion (1536d) & 30.5\% & First 30\%+ \\
10 & Vocabulary Normalization & 30.7\% & Minor gain \\
11 & MLP Neural Network & 28.8\% & Regression \\
12 & Egyptian$\to$German & 12.9\% & Exploratory \\
13 & Ridge Alpha Tuning & 31.6\% & Improved \\
14 & Iterative Procrustes + Hub Filter & 31.6\% & No change \\
15 & FastText Retraining (mc5\_w10) & \textbf{32.4\%} & \textbf{SOTA} \\
\bottomrule
\end{tabular}
\caption{Accuracy progression across 15 iterations.}
\label{tab:progression}
\end{table}

\subsection{Key Lessons}

\textbf{Linear methods beat neural networks.} With only $\sim$5,000 training pairs, neural architectures (MLPs, GANs, BERT) consistently overfit. The Ridge regression closed-form solution outperforms all learned mappings. This is the central methodological result: for low-resource alignment, reach for linear algebra before deep learning.

\textbf{BERT destroys ancient languages.} The WordPiece tokenizer fragments hieroglyphic transliterations into meaningless subwords: \textit{\d{h}r,w} (Horus) becomes \texttt{[\d{h}, \#\#r, \#\#,, \#\#w]}. Accuracy collapsed from 24.5\% to 0.5\%. Pre-trained NLP models encode modern linguistic assumptions that do not transfer to ancient languages.

\textbf{Quality beats quantity for anchors.} Adding 368 Coptic cognate anchors (V8) \textit{reduced} accuracy. Etymology $\neq$ semantics: words change meaning over millennia. A smaller set of high-confidence anchors outperforms a larger set of mixed-quality ones.

\textbf{MSE is anti-correlated with retrieval accuracy.} Cross-validated Ridge regression selects $\alpha = 100$ by mean squared error. But the actual best retrieval accuracy occurs at $\alpha = 0.001$---five orders of magnitude lower. Lower regularization produces ``sharper'' vectors that are noisier in an L2 sense but more discriminative for nearest-neighbor lookup. \textit{Optimizing for regression loss optimizes against retrieval.}

\textbf{Zero-padding helps via regularization.} Concatenating 768 zeros (placeholder visual features) to each 768-dimensional text vector improves accuracy despite adding no information. The larger parameter space ($1536 \times 300$ vs $768 \times 300$) acts as implicit regularization, preventing overfitting. This was confirmed by ablation: at matched regularization strength, 768d and 1536d perform identically.

\textbf{65\% of vocabulary is noise.} Over half of the Egyptian vocabulary consists of words appearing exactly once (hapax legomena). Their vectors are random. Filtering them (\texttt{min\_count=5}) reduces vocabulary by 87\% and \textit{increases} accuracy. The model was learning from noise it didn't need to learn from.

\textbf{60\% of the test set is function words.} The accuracy metric is depressed by function word ambiguity---many Egyptian function words (prepositions, pronouns) map to the same few English stopwords. On content words alone, the model performs substantially better. The ``hubness problem'' (82\% of predictions being ``the'') is not a model failure---it is the model correctly reflecting that function words dominate both languages.

% ══════════════════════════════════════════════════════════════════════════════
\section{Discoveries}
\label{sec:discoveries}

The following findings emerge from probing the aligned V15 embedding space with semantic analogies, concept midpoints, and cross-domain cluster analysis. They represent what literal translation cannot capture: the relational structure of how Egyptians connected ideas.

\subsection{Gold Is Divine Flesh \hfill \hiero{𓋴𓈖𓃀𓅱}}

The midpoint of ``gold'' and ``divine'' in English, projected into the Egyptian space, finds \textit{n\d{t}ri} (divine, similarity 0.642) and \textit{nbw} (gold, similarity 0.617) as the top results. They are geometrically colocated.

Modern readers treat the Egyptian statement ``the flesh of the gods is gold'' as poetic metaphor. The embedding space suggests otherwise. Gold and divinity are not \textit{compared} in the Egyptian textual record---they are \textit{colocated}. The same linguistic contexts that invoke divinity invoke gold, and vice versa. The model cannot distinguish them because the corpus does not distinguish them. This is not metaphor. It is ontology.

\subsection{Silence Is the Condition of the Dead \hfill \hiero{𓇯𓂋𓊽}}

The midpoint of ``silence'' and ``death'' returns variants of \textit{m(w)t} (die/dead) as all five nearest neighbors. There is no Egyptian word between silence and death---they occupy the same point in semantic space.

The Egyptians called the necropolis \textit{t\=a-sgr}---``the silent land.'' The dead were \textit{sgr.w}---``the silent ones.'' The embedding space confirms that this was not euphemism. Silence was death's defining quality. What the dead lost was not life in some abstract sense. It was voice.

\subsection{Seeing Was an Act of Magical Power \hfill \hiero{𓂀𓎛𓋴}}

The midpoint of ``eye'' and ``knowledge'' finds \textit{jr,t} (eyes) at the top, then \textit{\d{h}k\=a} (heka/magic) as the third-nearest word. Egyptian sight sits between knowing and spellcasting.

The Eye of Horus (\textit{w\underline{d}\=a,t}) was simultaneously an organ of perception, a protective amulet, and a unit of mathematical measurement. Seeing, in the Egyptian worldview, was not passive reception. It was an active, potent act---closer to divination than observation. The embedding space captures this without being told.

\subsection{The Snake Is Divine, Not Wise \hfill \hiero{𓆙𓊽𓋴}}

The midpoint of ``snake'' and ``wisdom'' finds gods and divinity---not wisdom, knowledge, or cunning. Every top result is a \textit{n\d{t}r} variant.

This directly contradicts the Greek-influenced reading of serpent symbolism. In the Greek tradition, the snake is associated with wisdom (the serpent of Asclepius, the caduceus). In the Egyptian tradition, the serpent---the uraeus cobra on the pharaoh's brow, the Apophis serpent of chaos---is associated with \textit{divine power}. The embedding space separates these cultural readings without any awareness of Greek mythology.

\subsection{Temple Is to House as God Is to Man \hfill \hiero{𓉐𓊽𓀀}}

The vector analogy \textit{house : temple :: man : ?} returns \textit{n\d{t}r} (god). The proportional relationship is preserved across the 4,000-year alignment: a temple is a god's house, in the same geometric sense that a house is a man's dwelling.

This is one of the few cases where vector arithmetic---the technique that produces ``king $-$ man $+$ woman $=$ queen'' in English---works across an ancient language boundary.

\subsection{Mother Is Royalty, Not Earth \hfill \hiero{𓅭𓏏𓀀}}

The midpoint of ``mother'' and ``earth'' does not find earth, soil, or land. It finds \textit{\d{h}m,t-nzw} (royal wife) and \textit{z\=a,t-nzw} (king's daughter). In the Egyptian textual record, motherhood is a \textit{royal} concept.

The ``earth mother''---the fertile, nurturing ground-goddess---is an Indo-European archetype (Gaia, Terra, Demeter). The Egyptian mother-goddess (Isis, Hathor) is not earthy. She is regal. The embedding space reveals this cultural divergence without being told about either tradition.

\subsection{Truth and Power Are the Same Force \hfill \hiero{𓁹𓋴𓊽}}

The midpoint of ``truth'' and ``power'' finds \textit{s\underline{h}m} (power/authority) at \#1 and \textit{\underline{h}ft.pl} (enemies) at \#4. Truth, power, and the defeat of enemies occupy the same semantic region.

This reflects \textit{m\=a\c{a}.t}---usually translated as ``truth'' or ``justice'' but encompassing cosmic order, righteous authority, and the subduing of chaos. Truth in the Egyptian worldview is not passive correctness. It is the active force that maintains the universe against its enemies.

\subsection{Love and Fear Meet at Eternity \hfill \hiero{𓆣𓇯𓋴}}

The midpoint of ``love'' and ``fear'' finds \textit{r-(n)\d{h}\d{h}} (eternal/forever) as the second-nearest word. Between love and fear, the Egyptians placed eternity.

This may reflect the theology of divine favor: the gods' love is awe-inspiring, their wrath is terrifying, and both extend beyond time. The offering formula---``that he may be loved and feared for eternity''---collapses these concepts into a single liturgical act.

% ══════════════════════════════════════════════════════════════════════════════
\section{Discussion}
\label{sec:discussion}

\subsection{Limitations}

The primary limitation is corpus bias. The surviving Egyptian texts are overwhelmingly funerary and religious. The embedding space reflects the language of temples, tombs, and offering formulas---not the language of markets, households, or daily conversation. Discoveries about ``the Egyptian worldview'' are more precisely about the \textit{literate priestly worldview preserved in stone}.

The 32.35\% Top-1 accuracy, while significant for unsupervised cross-lingual alignment of a dead language, means that two-thirds of test words do not retrieve their correct translation. The geometric relationships we analyze are \textit{statistical tendencies}, not certainties.

Additionally, the anchor pairs are extracted from German translations, introducing a second layer of cultural mediation. An Egyptian$\to$German$\to$English pipeline may introduce biases that a direct Egyptian$\to$English approach would avoid.

\subsection{What Translation Misses}

The core contribution of this work is not the translations---Egyptologists already have those---but the \textit{distances}. A dictionary can tell you that \textit{n\d{t}r} means ``god'' and \textit{nbw} means ``gold.'' But only the embedding space can tell you that they are essentially the same word, distributionally speaking. This is the kind of information that translation destroys: the proximity, the neighborhood, the geometric context.

Every act of translation is a projection from a high-dimensional space to a lower one. Meaning is lost. The embedding approach preserves more of the original dimensionality, allowing us to see relationships that the projection to English text made invisible.

\subsection{Generalizability}

The methodology is domain-agnostic. The question ``can we align two signal spaces that share latent structure?'' applies wherever:

\begin{itemize}[nosep]
    \item Two systems produce high-dimensional representations of related phenomena
    \item A small number of anchor correspondences can be identified
    \item The underlying structure is approximately isomorphic
\end{itemize}

Ancient Egyptian is a particularly compelling test case because it maximizes the difficulty: the languages are separated by 4,000 years, the corpus is small, and the conceptual categories are genuinely different. If vec2vec alignment works here, it can work in less extreme settings.

% ══════════════════════════════════════════════════════════════════════════════
\section{Conclusion}
\label{sec:conclusion}

We set out to build a lens---a way of reading the Ancient Egyptian textual record that reveals what literal translation leaves behind. The lens is imperfect. It resolves some features clearly and blurs others. But where it does resolve, it shows us a world organized differently from our own: a world where gold \textit{is} divinity, where silence \textit{is} death, where the snake is sacred rather than wise, where truth is a weapon and motherhood is a crown.

These discoveries emerged not from sophisticated neural architectures or massive compute, but from the simplest possible alignment: a linear projection learned from a few thousand anchor pairs. The central technical lesson is that for low-resource cross-lingual alignment, simplicity is not a compromise---it is the optimal strategy. The central humanistic lesson is that the geometry of a language encodes its worldview, and that worldview is recoverable even after 4,000 years of silence.

Translation gave us the words. The vectors gave us the world between them.

% ══════════════════════════════════════════════════════════════════════════════
\bibliographystyle{plainnat}
\begin{thebibliography}{99}

\bibitem[Bojanowski et~al., 2017]{bojanowski2017enriching}
Bojanowski, P., Grave, E., Joulin, A., and Mikolov, T. (2017).
\newblock Enriching word vectors with subword information.
\newblock \textit{Transactions of the Association for Computational Linguistics}, 5:135--146.

\bibitem[BBAW, 2023]{bbaw_tla}
Berlin-Brandenburg Academy of Sciences and Humanities (2023).
\newblock Thesaurus Linguae Aegyptiae (TLA): Digital corpus of Egyptian texts.
\newblock \url{https://aaew.bbaw.de/tla/}

\bibitem[Conneau et~al., 2018]{conneau2018word}
Conneau, A., Lample, G., Ranzato, M., Denoyer, L., and J\'{e}gou, H. (2018).
\newblock Word translation without parallel data.
\newblock In \textit{Proceedings of ICLR 2018}.

\bibitem[Firth, 1957]{firth1957synopsis}
Firth, J.~R. (1957).
\newblock A synopsis of linguistic theory, 1930--1955.
\newblock In \textit{Studies in Linguistic Analysis}, pages 1--32. Blackwell.

\bibitem[HamdiJr, 2023]{hamdijr_hieroglyphs}
HamdiJr (2023).
\newblock Egyptian hieroglyphs dataset.
\newblock \url{https://huggingface.co/datasets/HamdiJr/Egyptian_hieroglyphs}

\bibitem[Harris, 1954]{harris1954distributional}
Harris, Z.~S. (1954).
\newblock Distributional structure.
\newblock \textit{Word}, 10(2-3):146--162.

\bibitem[Lample et~al., 2018]{lample2018word}
Lample, G., Conneau, A., Denoyer, L., and Ranzato, M. (2018).
\newblock Unsupervised machine translation using monolingual corpora only.
\newblock In \textit{Proceedings of ICLR 2018}.

\bibitem[Mikolov et~al., 2013a]{mikolov2013efficient}
Mikolov, T., Chen, K., Corrado, G., and Dean, J. (2013a).
\newblock Efficient estimation of word representations in vector space.
\newblock In \textit{Proceedings of ICLR 2013 Workshop}.

\bibitem[Mikolov et~al., 2013b]{mikolov2013linguistic}
Mikolov, T., Yih, W., and Zweig, G. (2013b).
\newblock Linguistic regularities in continuous space word representations.
\newblock In \textit{Proceedings of NAACL-HLT 2013}, pages 746--751.

\bibitem[Pennington et~al., 2014]{pennington2014glove}
Pennington, J., Socher, R., and Manning, C.~D. (2014).
\newblock {GloVe}: Global vectors for word representation.
\newblock In \textit{Proceedings of EMNLP 2014}, pages 1532--1543.

\bibitem[Sch\"{o}nemann, 1966]{schonemann1966generalized}
Sch\"{o}nemann, P.~H. (1966).
\newblock A generalized solution of the orthogonal {P}rocrustes problem.
\newblock \textit{Psychometrika}, 31(1):1--10.

\end{thebibliography}

\end{document}
